\documentclass[11pt,a4paper]{article}

\usepackage[T1]{fontenc}
\usepackage{newtxtext,newtxmath}
\usepackage[margin=25mm,headheight=15pt]{geometry}
\usepackage{microtype}
\usepackage{amsmath,mathtools,bm}
\usepackage{booktabs,tabularx,array,longtable,multirow}
\usepackage{graphicx,float,subcaption,pdflscape}
\usepackage[dvipsnames,table]{xcolor}
\usepackage{tikz}
\usetikzlibrary{arrows.meta,positioning,fit,backgrounds,calc}
\usepackage[most]{tcolorbox}
\usepackage{enumitem}
\usepackage{caption}
\usepackage{fancyhdr}
\usepackage{titlesec}
\usepackage{placeins}
\usepackage{hyperref}
\usepackage[nameinlink,noabbrev]{cleveref}

\definecolor{Navy}{HTML}{17324D}
\definecolor{Blue}{HTML}{2F6B9A}
\definecolor{Teal}{HTML}{2A9D8F}
\definecolor{Gold}{HTML}{E9C46A}
\definecolor{Orange}{HTML}{F4A261}
\definecolor{Red}{HTML}{D1495B}
\definecolor{Light}{HTML}{E8EEF3}
\definecolor{Ink}{HTML}{263238}
\definecolor{Muted}{HTML}{66717D}

\hypersetup{
  colorlinks=true,
  linkcolor=Navy,
  citecolor=Blue,
  urlcolor=Blue,
  pdftitle={SPX Volatility-Surface Risk on WRDS, 2005-2021},
  pdfauthor={Volatility Surface Risk Lab},
  pdfsubject={Reproducible thesis project report},
  pdfkeywords={SPX, implied volatility, VaR, PCA, Heston, PSP, WRDS, OptionMetrics}
}

\pagestyle{fancy}
\fancyhf{}
\fancyhead[L]{\small\color{Muted}\nouppercase{\leftmark}}
\fancyhead[R]{\small\color{Muted}SPX volatility-surface risk}
\fancyfoot[L]{\small\color{Muted}WRDS OptionMetrics | 2005-2021}
\fancyfoot[R]{\small\color{Navy}\thepage}
\renewcommand{\headrulewidth}{0.4pt}
\renewcommand{\footrulewidth}{0pt}

\titleformat{\section}
  {\Large\bfseries\color{Navy}}{\thesection}{0.65em}{}
\titleformat{\subsection}
  {\large\bfseries\color{Blue}}{\thesubsection}{0.65em}{}
\titleformat{\subsubsection}
  {\normalsize\bfseries\color{Ink}}{\thesubsubsection}{0.65em}{}
\titlespacing*{\section}{0pt}{2.1ex plus .6ex}{1.0ex}
\titlespacing*{\subsection}{0pt}{1.8ex plus .5ex}{0.8ex}

\captionsetup{font=small,labelfont={bf,color=Navy},textfont={color=Ink}}
\setlist{nosep,leftmargin=1.5em}
\setlength{\parindent}{0pt}
\setlength{\parskip}{5.5pt plus 1.5pt}
\setcounter{tocdepth}{2}
\setcounter{secnumdepth}{3}
\graphicspath{{../long_horizon_2005_2021/figures/}}

\newcolumntype{Y}{>{\raggedright\arraybackslash}X}
\newcolumntype{C}{>{\centering\arraybackslash}X}
\newcommand{\yes}{\textcolor{Teal}{\textbf{Pass}}}
\newcommand{\no}{\textcolor{Red}{\textbf{Reject}}}
\newcommand{\code}[1]{\texttt{#1}}
\newcommand{\RawRows}{8,671,456}
\newcommand{\EligibleRows}{6,712,903}
\newcommand{\SurfaceDates}{4,280}
\newcommand{\ExtrapolatedShare}{24.00\%}
\newcommand{\PCAExplained}{96.71\%}
\newcommand{\CorrelationPCOneCosine}{0.994}
\newcommand{\CorrelationPCTwoCosine}{0.898}
\newcommand{\CorrelationPCThreeCosine}{0.849}
\newcommand{\CorrelationPCFiveCosine}{0.648}
\newcommand{\HestonMedianRMSE}{1.10\%}
\newcommand{\HestonPninetyRMSE}{1.90\%}
\newcommand{\HestonFellerShare}{3.1\%}
\newcommand{\HestonXiBoundShare}{32.7\%}
\newcommand{\HestonNinetyBreaches}{371}
\newcommand{\HestonNinetyFiveBreaches}{186}
\newcommand{\HestonNinetyNineBreaches}{36}
\newcommand{\FullGBMNinetyBreaches}{382}
\newcommand{\FullGBMNinetyFiveBreaches}{187}
\newcommand{\FullGBMNinetyNineBreaches}{37}
\newcommand{\FullHestonPNinetyBreaches}{333}
\newcommand{\FullHestonPNinetyFiveBreaches}{200}
\newcommand{\FullHestonPNinetyNineBreaches}{72}
\newcommand{\FullHestonPNinetyFiveCCP}{0.947}
\newcommand{\FullPairedCorrelation}{-0.659}
\newcommand{\FullExtrapolatedShare}{29.0\%}
\newcommand{\CurveFallbackCount}{1}
\newcommand{\PhysicalKappaBoundShare}{7.7\%}
\newcommand{\HestonMCMeanError}{0.034\%}
\newcommand{\HestonPriceGridError}{3.38}


\newtcolorbox{keybox}[1][]{
  enhanced,breakable,colback=Light!55,colframe=Blue,
  boxrule=0.7pt,arc=2pt,left=8pt,right=8pt,top=7pt,bottom=7pt,
  fonttitle=\bfseries,title=#1
}
\newtcolorbox{warningbox}[1][]{
  enhanced,breakable,colback=Gold!18,colframe=Orange!85!black,
  boxrule=0.7pt,arc=2pt,left=8pt,right=8pt,top=7pt,bottom=7pt,
  fonttitle=\bfseries,title=#1
}

\begin{document}
\pagenumbering{roman}

\begin{titlepage}
\begin{tikzpicture}[remember picture,overlay]
  \fill[Navy] (current page.north west) rectangle ([yshift=-10.4cm]current page.north east);
  \fill[Blue] (current page.south west) rectangle ([xshift=12mm]current page.north west);
  \fill[Teal] ([xshift=12mm,yshift=22mm]current page.south west)
    rectangle ([xshift=19mm,yshift=83mm]current page.south west);
  \draw[Gold,line width=2pt] ([xshift=25mm,yshift=-57mm]current page.north west)
    -- ([xshift=-25mm,yshift=-57mm]current page.north east);
\end{tikzpicture}
\vspace*{16mm}
{\color{white}\sffamily\bfseries\fontsize{15}{18}\selectfont MASTER THESIS PROJECT REPORT\par}
\vspace{16mm}
{\color{white}\sffamily\bfseries\fontsize{29}{34}\selectfont
SPX Volatility-Surface Risk\par}
\vspace{3mm}
{\color{Gold}\sffamily\bfseries\fontsize{19}{23}\selectfont
WRDS OptionMetrics, 2005-2021\par}

\vspace{30mm}
{\color{Navy}\sffamily\bfseries\fontsize{17}{21}\selectfont
Surface reconstruction, PCA, Heston, and VaR validation\par}
\vspace{10mm}
{\color{Ink}\large
Theory, implementation, computational design, empirical evidence, and audit trail\par}

\vfill
\begin{tabular}{@{}ll@{}}
\textbf{Project} & Volatility Surface Risk Lab \\
\textbf{Analysis window} & 3 January 2005 to 31 December 2021 \\
\textbf{Data source} & WRDS OptionMetrics IvyDB US \\
\textbf{Prepared} & 18 August 2026 \\
\textbf{Version} & Reproducible research release 4.0
\end{tabular}
\vspace{14mm}
\end{titlepage}
\setcounter{page}{2}

\section*{Executive summary}
\addcontentsline{toc}{section}{Executive summary}

This report documents the consolidation and extension of the SPX option-risk thesis project. The
production analysis uses \RawRows\ WRDS option records, retains \EligibleRows\ observations after
quality controls, and reconstructs \SurfaceDates\ daily implied-volatility surfaces. It evaluates the
original PSP, covariance-PCA, and PWG engines plus correlation-PCA and a daily calibrated Heston
surface-history model at 90\%, 95\%, and 99\% VaR. A second, deliberately separate experiment holds
option strikes fixed, simulates spot and variance jointly under Heston, and fully reprices a daily
delta-hedged call strip. The code, configuration, machine-readable tables, and figures are generated
from one reproducible Python pipeline.

The leading five covariance principal components explain \PCAExplained\ of daily surface-change
variation, but the original loading figure was not a clean level--slope--curvature decomposition.
Node standard deviations differ by more than a factor of twelve, so the raw covariance basis is
dominated by the short-dated wing. Correlation-PCA recovers a level-like PC1, moneyness-slope PC2,
term-slope PC3, and smile-curvature PC5. PSP, both PCA variants, and Heston-HS pass unconditional,
independence, and joint conditional-coverage tests at all three levels. PWG fails joint coverage at
all levels.

Those conclusions belong to the controlled vega-only experiment. In the nonlinear full-revaluation
experiment, physical-measure Heston passes joint conditional coverage only at 95\%; the result does
not support a blanket claim that Heston solves the tail-risk problem.

\begin{keybox}[Central interpretation]
The strongest result is not that one model wins every criterion. In the surface-shock experiment,
Heston-HS is competitive but does not significantly dominate PSP. In the full-revaluation
experiment, Heston-MC-P passes joint conditional coverage at 95\%
($p=\FullHestonPNinetyFiveCCP$), but rejects at 90\% and 99\%. The paired GBM+PSP-FHS benchmark is
closest to nominal unconditional coverage at all three levels, although its 95\% and 99\% exceptions
cluster. Heston-MC-Q is a pricing-measure sensitivity, not a real-world VaR model, and fails joint
coverage at every level.
\end{keybox}

\begin{warningbox}[Scope boundary]
Two distinct P\&L experiments are reported and never pooled. The first isolates implied-volatility
risk with $\Delta V\approx\bm{w}^{\top}\Delta\bm{\sigma}$. The second uses fixed strikes, maturity
roll-down, a daily delta hedge, financing and dividends, joint spot--variance scenarios, and nonlinear
option repricing. Its Heston-P transition is an estimated physical proxy while its conditional option
prices remain risk-neutral. Neither experiment is described as a live trading portfolio.
\end{warningbox}

\clearpage
\tableofcontents
\clearpage
\listoffigures
\clearpage
\listoftables
\clearpage
\pagenumbering{arabic}

\section{Research objective and scope}

The project estimates the market risk induced by movements of the SPX implied-volatility surface.
It replaces distributed OptionsData files and ad hoc notebooks with date-bounded WRDS extraction,
auditable filters, fixed-grid surface reconstruction, rolling factor analysis, scenario generation,
and formal out-of-sample forecast evaluation. The requested endpoint is 31 December 2021; no 2022
or later observations enter the results.

The maintained analysis contains a controlled surface-risk experiment and a separate full-revaluation
experiment. The surface-risk methods are:

\begin{itemize}
  \item \textbf{PSP}: recency-weighted historical projection of full volatility-surface shocks;
  \item \textbf{PCA}: rolling factor compression and bootstrap of PCA-reconstructed shocks;
  \item \textbf{PWG}: a full-grid multivariate Gaussian benchmark with Ledoit-Wolf shrinkage.
  \item \textbf{PCA-C}: the same rolling PCA risk engine after node-wise standardization, included
  to test level--slope--curvature interpretation without changing the observed P\&L.
  \item \textbf{Heston-HS}: historical shocks of a daily calibrated Heston-implied-volatility
  surface, expressed on the same 45-node grid.
\end{itemize}

The full-revaluation experiment compares:
\begin{itemize}
  \item \textbf{GBM+PSP-FHS}: a geometric spot update driven by the historical standardized return
  innovation paired with the same day's complete surface shock;
  \item \textbf{Heston-MC-P}: joint spot--variance simulation under trailing physical parameters,
  followed by conditional risk-neutral Heston repricing; and
  \item \textbf{Heston-MC-Q}: the same simulation under calibrated risk-neutral dynamics, retained
  only as a pricing-measure sensitivity check.
\end{itemize}

The legacy PwG interest-rate folder is excluded. The zero curve is taken only from WRDS
OptionMetrics \code{zerocd}; the dividend yield is taken from \code{idxdvd}. The name PWG is retained
only for the full-grid Gaussian volatility benchmark.

\section{End-to-end model pipeline}

\begin{figure}[H]
\centering
\includegraphics[width=\textwidth]{00_pipeline.png}
\caption{Reproducible model pipeline. The upper branch
isolates surface risk; the lower branch performs fixed-strike full revaluation under paired PSP and
joint Heston scenarios. Both use the same dates, grid, tail levels, and validation layer. Every
forecast uses information available at or before its origin.}
\label{fig:pipeline}
\end{figure}

The pipeline separates data engineering, surface estimation, risk generation, and validation. This
separation is essential: a visually smooth surface is not automatically a good risk forecast, and a
correct exception count does not guarantee independent exceptions or low forecast loss.

\section{Data and preprocessing}

\subsection{WRDS mapping}

SPX identifiers are discovered from the OptionMetrics security-name history rather than hard-coded.
The annual option table supplies quote date, expiration, call/put flag, scaled strike, bid, offer,
volume, open interest, implied volatility, Greeks, and vendor forward. The annual security-price
table supplies the underlying close; the zero curve is linearly interpolated by days to expiration.
All rates are normalized to decimal form before use.

For maturity $T$ in years, the model fallback forward is
\begin{equation}
F_{t,T}=S_t\exp\!\left[(r_{t,T}-q_{t,T})T\right],
\end{equation}
and forward moneyness and normalized moneyness are
\begin{equation}
m=\frac{F_{t,T}}{K}, \qquad x=\frac{\log(F_{t,T}/K)}{\sqrt{T}}.
\end{equation}

\subsection{Production filters}

The long-horizon run uses calls, mid-price at least \$1, 15--550 days to expiration, forward
moneyness 0.70--1.30, relative spread no greater than 1.0, nonnegative open interest, positive
implied volatility, and at least 25 eligible observations per day. The fitted reporting grid covers
moneyness 0.80--1.20 and maturities 30, 60, 90, 180, and 360 days.

\begin{table}[H]
\centering
\caption{Verified data and analysis dimensions}
\label{tab:data-summary}
\begin{tabular}{lr}
\toprule
Quantity & Verified value \\
\midrule
WRDS option rows returned after server-side restrictions & \RawRows \\
Eligible option observations after research filters & \EligibleRows \\
Daily reconstructed surfaces & \SurfaceDates \\
Grid nodes per surface & 45 \\
Out-of-hull grid-node share & \ExtrapolatedShare \\
Out-of-sample VaR observations per model and level & 4,027 \\
Scenario draws per forecast and model & 2,000 \\
Rolling window / minimum history & 500 / 252 days \\
\bottomrule
\end{tabular}
\end{table}

\begin{figure}[H]
\centering
\includegraphics[width=\textwidth]{01_data_coverage.png}
\caption{Eligible SPX option coverage through time after all production filters.}
\label{fig:data-coverage}
\end{figure}

\section{Black-Scholes-Merton and volatility-surface theory}

\subsection{European option prices}

Under geometric Brownian motion with constant volatility, frictionless trading, continuous
rebalancing, constant continuously compounded $r$ and dividend yield $q$, the European call and put
prices are
\begin{align}
C(S,K,T,r,q,\sigma)&=Se^{-qT}N(d_1)-Ke^{-rT}N(d_2),\\
P(S,K,T,r,q,\sigma)&=Ke^{-rT}N(-d_2)-Se^{-qT}N(-d_1),\\
d_1&=\frac{\log(S/K)+(r-q+\tfrac12\sigma^2)T}{\sigma\sqrt{T}},
\qquad d_2=d_1-\sigma\sqrt{T}.
\end{align}
Here $N(\cdot)$ and $\phi(\cdot)$ denote the standard-normal distribution and density. The new
\code{black\_scholes.py} module implements calls, puts, delta, gamma, vega, theta, rho, and bracketed
Brent inversion for implied volatility, with unit tests against standard benchmark values.

\subsection{Implied volatility and Greeks}

Implied volatility $\widehat\sigma$ is the root of
\begin{equation}
BS(S,K,T,r,q,\widehat\sigma)-P^{\mathrm{mkt}}=0.
\end{equation}
The production run uses the OptionMetrics standardized implied volatility when it is positive. If
\code{filters.impute\_missing\_iv=true}, the pipeline instead inverts the mid-price through the tested
Black-Scholes-Merton routine. This switch is off in the reported run, so no imputed IV silently enters
the empirical sample.

The first-order sensitivity to implied volatility is
\begin{equation}
\mathrm{Vega}=\frac{\partial V}{\partial\sigma}=Se^{-qT}\phi(d_1)\sqrt{T}.
\end{equation}
Other relevant sensitivities are delta $\partial V/\partial S$, gamma
$\partial^2V/\partial S^2$, theta $\partial V/\partial t$, and rho $\partial V/\partial r$. A local
price approximation is
\begin{equation}
\Delta V\approx \Delta\,\Delta S+\tfrac12\Gamma(\Delta S)^2+\mathrm{Vega}\,\Delta\sigma
+\mathrm{Theta}\,\Delta t+\mathrm{Rho}\,\Delta r.
\end{equation}
The first experiment sets the non-vega terms aside to identify volatility-surface risk cleanly. On
the $5\times9$ grid, a disclosed positive exposure vector is normalized to total vega 500,000:
\begin{equation}
w_{T,m}\propto \sqrt{T}\exp\!\left[-\frac12\left(\frac{m-1}{0.14}\right)^2\right],
\qquad \sum_{T,m}w_{T,m}=500{,}000.
\end{equation}

\begin{table}[H]
\centering
\caption{Vega exposure by maturity; node minima and maxima reflect the smile weighting}
\label{tab:vega}
\begin{tabular}{rrrrr}
\toprule
Maturity (days) & Total vega & Share & Minimum node & Maximum node \\
\midrule
30 & 49,703 & 9.9\% & 2,856 & 7,923 \\
60 & 70,290 & 14.1\% & 4,039 & 11,204 \\
90 & 86,087 & 17.2\% & 4,946 & 13,722 \\
180 & 121,746 & 24.3\% & 6,995 & 19,406 \\
360 & 172,175 & 34.4\% & 9,892 & 27,444 \\
\bottomrule

\end{tabular}
\end{table}

The full-revaluation experiment uses this disclosed profile only as a quantity shape. One common
daily scalar makes the call strip's actual Black--Scholes vega equal to 500,000; it does not divide
node by node by very small wing vegas. Spot, convexity, time decay, financing, dividends, and the
nonlinear volatility effect then enter through complete option-price differences.

\subsection{Why a surface is necessary}

Market option prices violate the constant-volatility implication: implied volatility varies with
strike/moneyness and maturity, forming a skewed surface. A single volatility number discards the
relative movement of short versus long maturities and wings versus at-the-money contracts. The
research object is therefore $\sigma_t(m,T)$ and especially its daily change
$\Delta\sigma_t(m,T)$.

\subsection{Heston stochastic volatility and its implied surface}

Black--Scholes does not need to be the true market dynamics for implied volatility to be a useful
coordinate. Heston replaces constant variance with a correlated square-root process
\cite{heston1993}:
\begin{align}
\frac{dS_t}{S_t}&=(r-q)\,dt+\sqrt{v_t}\,dW_t^S,\\
dv_t&=\kappa(\theta-v_t)\,dt+\xi\sqrt{v_t}\,dW_t^v,
\qquad d\langle W^S,W^v\rangle_t=\rho\,dt.
\end{align}
Here $v_0$ is current variance, $\kappa$ mean reversion, $\theta$ long-run variance, $\xi$ the
volatility of variance, and $\rho$ the spot--variance correlation. Heston supplies option prices,
not a directly observed ``Heston volatility surface.'' The comparable surface is obtained by
converting each Heston price back to its Black--Scholes-equivalent implied volatility.

The displayed drift $r-q$ and parameters inferred from option prices belong to the risk-neutral
measure $Q$. VaR forecasts concern the physical distribution $P$. The full-revaluation experiment
therefore keeps these roles separate: trailing return and variance-state histories estimate the
$P$ transition, while conditional option values at the horizon use the $Q$ pricing function. A
simulation that uses $Q$ for both is reported only to show the consequence of confusing those roles.

The empirical grid is already in forward moneyness $m=F/K$. The calibration therefore normalizes
$F=1$, sets $K=1/m$, and removes the common discount factor. For each date, the code minimizes
vega-scaled price errors
\begin{equation}
\widehat{\bm\psi}_t=\operatorname*{arg\,min}_{\bm\psi}
\sum_{i=1}^{45}\left[
\frac{C_H(1,K_i,T_i;\bm\psi)-
C_{BS}(1,K_i,T_i;\sigma^{\mathrm{mkt}}_{t,i})}
{\max(\mathrm{Vega}_{t,i},0.025)}
\right]^2+\text{light warm-start regularization},
\end{equation}
where $\bm\psi=(v_0,\kappa,\theta,\xi,\rho)$. Calls are evaluated by damped characteristic-function
inversion with Gauss--Laguerre quadrature, following the transform approach of
\cite{carrmadan1999}; the recovered surface is
\begin{equation*}
\sigma^H_{t,i}=\operatorname{BS}^{-1}\bigl(C_H(1,K_i,T_i;\widehat{\bm\psi}_t)\bigr).
\end{equation*}
All five parameters are warm-started and recalibrated every day. This is slower than a monthly
filter but prevents stale skew and term parameters from masquerading as a Heston failure. The
Feller margin $2\kappa\theta-\xi^2$ and parameter-bound hits are reported as diagnostics rather than
silently imposed constraints.

\section{Volatility-surface reconstruction}

All four estimators model $y=\log\sigma$ on normalized moneyness $x$ and maturity $T$. The log target
enforces positive fitted volatility before the explicit 3\%--200\% bounds.

\subsection{Methods}

\begin{enumerate}
  \item \textbf{Polynomial / Dumas form:}
  \begin{equation}
  \log\sigma=a_0+a_1x+a_2T+a_3x^2+a_4xT+\varepsilon.
  \end{equation}
  It is interpretable and fast but can underfit local smile curvature.

  \item \textbf{Tensor B-spline:} cubic basis functions $B_j(x)$ and $C_k(T)$ are crossed and
  ridge-regularized,
  \begin{equation}
  \widehat{\bm\beta}=\arg\min_{\bm\beta}
  \|\bm y-\bm X\bm\beta\|_2^2+\lambda\|\bm\beta\|_2^2,
  \qquad X_{i,jk}=B_j(x_i)C_k(T_i).
  \end{equation}
  This is the production surface.

  \item \textbf{Thin-plate radial basis:} scaled radial terms
  $\varphi(r)=r^2\log r$ are centered at quantile-selected knots and ridge-stabilized.

  \item \textbf{Linear Delaunay interpolation:} local planar interpolation is used inside the
  observed cloud, with nearest fitted values outside it.
\end{enumerate}

Every 21st trading day is evaluated with a deterministic date-wise 80/20 holdout. The fixed grid is
also tested against the daily convex hull in $(x,T)$ space. Nodes outside the quote cloud use the
nearest fitted observation and carry an \code{extrapolated=true} audit flag. This guard was crucial:
without it, unsupported spline tails can form impressive but artificial volatility spikes.

\begin{table}[H]
\centering
\caption{Held-out surface reconstruction performance over 221 validation dates}
\label{tab:surface-methods}
\begin{tabular}{lrrrrr}
\toprule
Method & Dates & Median IV RMSE & Mean IV RMSE & Median log-IV MAE & Median $R^2$ \\
\midrule
linear & 221 & 0.00367 & 0.00437 & 0.00574 & 0.9974 \\
b\_spline & 221 & 0.00370 & 0.00419 & 0.00924 & 0.9967 \\
thin\_plate & 221 & 0.00553 & 0.00596 & 0.01447 & 0.9924 \\
polynomial & 221 & 0.01209 & 0.01247 & 0.04073 & 0.9611 \\
\bottomrule

\end{tabular}
\end{table}

Linear interpolation has the lowest median IV RMSE by a small margin; the B-spline has the lower mean
RMSE and produces a stable dense surface. The production choice therefore balances predictive error,
smoothness, and downstream PCA compatibility rather than selecting by in-sample fit.

\begin{figure}[H]
\centering
\includegraphics[width=.92\textwidth]{12_surface_method_metrics.png}
\caption{Held-out reconstruction error and validation $R^2$ by method.}
\label{fig:surface-metrics}
\end{figure}

\begin{figure}[H]
\centering
\includegraphics[width=.90\textwidth]{11_surface_methods_3d.png}
\caption{Three-dimensional comparison of the four surface reconstruction families on the same date
and axes.}
\label{fig:surface-methods-3d}
\end{figure}

\begin{figure}[H]
\centering
\begin{subfigure}[t]{.49\textwidth}
  \includegraphics[width=\textwidth]{02_surface_snapshot.png}
  \caption{Heat map and maturity smile slices.}
\end{subfigure}\hfill
\begin{subfigure}[t]{.49\textwidth}
  \includegraphics[width=\textwidth]{08_reconstructed_surface_3d.png}
  \caption{Production B-spline surface in three dimensions.}
\end{subfigure}
\caption{Latest reconstructed SPX implied-volatility surface.}
\label{fig:surface-latest}
\end{figure}

\section{PCA and temporal risk engines}

\subsection{Principal component representation}

Let $\bm s_t\in\mathbb{R}^{45}$ be the fixed-grid surface and
$\Delta\bm s_t=\bm s_t-\bm s_{t-1}$. PCA diagonalizes the historical covariance matrix,
\begin{equation}
\widehat\Sigma=\bm V\bm\Lambda\bm V^{\top}, \qquad
\Delta\bm s_t\approx\sum_{j=1}^{k}z_{t,j}\bm v_j.
\end{equation}
The rolling risk model uses $k=5$. PCA is re-estimated at every forecast origin with the randomized
solver, so no future loading enters an earlier forecast.

\begin{table}[H]
\centering
\caption{Variance explained by the production PCA factors}
\label{tab:pca}
\begin{tabular}{lrr}
\toprule
Component & Individual share & Cumulative share \\
\midrule
PC1 & 77.03\% & 77.03\% \\
PC2 & 14.68\% & 91.71\% \\
PC3 & 2.48\% & 94.19\% \\
PC4 & 1.52\% & 95.71\% \\
PC5 & 0.99\% & 96.71\% \\
\bottomrule

\end{tabular}
\end{table}

\begin{figure}[H]
\centering
\includegraphics[width=.82\textwidth]{04_pca_explained_variance.png}
\caption{Individual and cumulative variance explained by the five PCA factors.}
\label{fig:pca-variance}
\end{figure}

\begin{figure}[H]
\centering
\includegraphics[width=\textwidth]{09_pca_components_3d.png}
\caption{Original covariance-PCA loading surfaces. These are empirical eigenvectors, not a clean
textbook level--slope--curvature basis: high-variance short-dated wing nodes dominate the geometry.}
\label{fig:pca-3d}
\end{figure}

The initial caption was therefore too confident. A level--slope--curvature interpretation is a common
heuristic, not a PCA theorem. In this sample, node standard deviations range from 0.59 to 7.45 IV
percentage points per day; the largest values occur at the short-dated high-$F/K$ wing. Covariance
PCA correctly maximizes variance in IV units, but its first vector need not be a parallel shift.

As a robustness check, correlation-PCA standardizes every node before eigendecomposition and compares
the first ten eigenvectors with orthonormal level, moneyness-slope, term-slope, and smile-curvature
templates. Signs are irrelevant, so alignment is reported as an absolute cosine.

\begin{table}[H]
\centering
\caption{Correlation-PCA alignment with canonical surface templates}
\label{tab:pca-interpretation}
\begin{tabular}{lrrr}
\toprule
Component & Variance share & Closest template & Absolute cosine \\
\midrule
PC1 & 79.94\% & Level & 0.994 \\
PC2 & 9.37\% & Moneyness slope & 0.898 \\
PC3 & 3.42\% & Term slope & 0.849 \\
PC5 & 1.08\% & Smile curvature & 0.648 \\
\bottomrule

\end{tabular}
\end{table}

PC1 is a standardized-node level shift ($\CorrelationPCOneCosine$), and PC2 is a moneyness
slope/skew ($\CorrelationPCTwoCosine$). PC3 is primarily a \emph{term} slope
($\CorrelationPCThreeCosine$), while smile convexity appears most clearly in PC5
($\CorrelationPCFiveCosine$). Thus the user's expected geometry does emerge after scaling, but the
two-dimensional maturity axis places term slope before strike curvature.

\begin{figure}[H]
\centering
\includegraphics[width=\textwidth]{18_pca_standardization_3d.png}
\caption{Covariance versus correlation PCA in three dimensions. Correlation loadings are shown in
standardized-node units; PC5 is included because it, rather than PC3, carries the clearest smile
curvature in this two-dimensional surface.}
\label{fig:pca-standardized}
\end{figure}

\begin{figure}[H]
\centering
\includegraphics[width=\textwidth]{19_pca_diagnostic_dashboard.png}
\caption{Canonical alignment, reconstruction trade-off, and subperiod stability. Five covariance
components reconstruct physical IV changes with 18.15\% relative error versus 26.39\% for
correlation PCA; the first three standardized directions are more stable than the full five-factor
subspace.}
\label{fig:pca-diagnostics}
\end{figure}

\begin{figure}[H]
\centering
\includegraphics[width=\textwidth]{10_pca_reconstruction_3d.png}
\caption{Observed latest surface, three-factor reconstruction, and residual surface.}
\label{fig:pca-reconstruction}
\end{figure}

\subsection{PSP, PCA, PWG, and Heston surface-history scenarios}

For a forecast origin $t$, all models use only changes observed through $t$; the $t+1$ realization
is excluded.

\textbf{PSP.} Historical full-grid shocks are sampled with replacement. With decay
$\lambda=0.995$, the probability assigned to an observation of age $a$ is
\begin{equation}
p_a=\frac{\lambda^a}{\sum_j\lambda^j}.
\end{equation}
This preserves empirical skewness, cross-node dependence, and stress shapes.

\textbf{PCA.} The historical window is projected into the five-factor space, reconstructed, mapped
to portfolio P\&L, and bootstrapped. This deliberately removes the low-variance residual surface.

\textbf{PWG.} The mean $\widehat{\bm\mu}$ and Ledoit-Wolf covariance
$\widehat\Sigma_{LW}$ imply the scalar Gaussian portfolio law
\begin{equation}
\Pi\sim N\!\left(\bm w^{\top}\widehat{\bm\mu},
\bm w^{\top}\widehat\Sigma_{LW}\bm w\right).
\end{equation}
Shrinkage stabilizes the 45-dimensional covariance, but normality still constrains tail shape.

\textbf{Heston-HS.} Each observed date is mapped to a daily calibrated Heston-implied surface
$\bm s_t^H$. The risk engine applies the same exponentially weighted 500-day historical simulation
to $\Delta\bm s_t^H$ and maps those structural surface shocks through the unchanged vega vector.
Realized P\&L remains the observed B-spline surface P\&L; Heston is used only to filter scenario
shapes. This is the cleanest like-for-like answer to how a stochastic-volatility model incorporates
the volatility surface.

\subsection{Joint Heston scenarios and fixed-strike full revaluation}

Heston-HS answers a surface-filtering question, not a joint market-dynamics question. The second
experiment therefore simulates both $S$ and $v$ and reprices the same call strip one trading day
later. At each origin, strikes are fixed at $K_i=F_{t,T_i}/m_i$. If $g_i$ is the normalized vega
shape in \cref{tab:vega}, the quantity and initial hedge are
\begin{equation}
n_{t,i}=a_t g_i,\qquad
a_t=\frac{500{,}000}{\sum_i g_i\,\mathrm{Vega}_{t,i}},\qquad
h_t=-\sum_i n_{t,i}\Delta_{t,i}.
\end{equation}
One common scalar avoids economically meaningless quantities at low-vega wing nodes. The cash
account makes the initial position zero-cost,
$B_t=-(\sum_i n_{t,i}C_{t,i}+h_tS_t)$. The realized one-day P\&L is
\begin{equation}
\Pi_{t+1}=\sum_i n_{t,i}C_{t+1,i}+h_tS_{t+1}
+B_t e^{r_{t,1}\Delta\tau}+h_tS_t\!\left(e^{q_t\Delta\tau}-1\right),
\end{equation}
where maturities roll down by the actual calendar gap. The horizon market price uses the next WRDS
curve and the next reconstructed IV surface evaluated at the new forward moneyness. Thus the
evaluation target includes delta-hedging error, gamma, vega, theta, dividends, and financing.

Daily option calibration identifies $Q$ parameters. To obtain a physical transition, a trailing
500-day regression uses the observable option-implied $v_0$ sequence as a variance-state proxy:
\begin{align}
\Delta v_{u+1}&=a+bv_u+\varepsilon_{u+1},\\
\widehat\kappa_P&=-b/\Delta t,\qquad
\widehat\theta_P=a/(\widehat\kappa_P\Delta t),\qquad
\widehat\xi_P=\operatorname{sd}\!\left(
\frac{\varepsilon_{u+1}}{\sqrt{v_u\Delta t}}\right),
\end{align}
and $\rho_P$ is the correlation between standardized spot and variance innovations. The total-return
drift $\mu_P$ is the winsorized mean close-to-close return plus dividend yield. Bounds are applied
only for numerical stability and every contact is reported. This is a transparent physical proxy,
not a claim that risk-neutral instantaneous variance is a perfectly observed physical state.

Heston-MC-P uses the physical estimates above for the one-day state transition. Heston-MC-Q instead
uses the current short rate and calibrated $Q$ parameters. Both use eight projected-Euler substeps:
\begin{align}
v_j^+&=\max(v_j,0),\qquad
Z_j^S=\rho Z_j^v+\sqrt{1-\rho^2}\,Z_j^\perp,\\
S_{j+1}&=S_j\exp\!\left[
(\mu-q-\tfrac12v_j^+)\delta+\sqrt{v_j^+\delta}\,Z_j^S\right],\\
v_{j+1}&=\max\!\left\{v_j+\kappa(\theta-v_j^+)\delta
+\xi\sqrt{v_j^+\delta}\,Z_j^v,0\right\}.
\end{align}
The projected scheme is explicit and biased at finite step size, so 4/8/16-step convergence is
reported rather than assumed \cite{lord2010}. At the horizon, options are priced with the current
$Q$ Heston characteristic function conditional on $(S_{t+1},v_{t+1})$. The current market-minus-
Heston price residual is frozen for one day and added back before no-arbitrage clipping. This
preserves today's price level while allowing state-dependent Heston curvature.

The comparator is filtered historical simulation, not independent mixing. A historical standardized
spot innovation and the entire PSP surface shock are sampled from the \emph{same} date with
exponential weights. The current variance rescales the spot innovation, the shocked surface is
evaluated at each simulated forward moneyness, and Black--Scholes fully reprices the fixed strikes.
The median trailing correlation between the paired spot innovation and the mean surface shock is
\FullPairedCorrelation, confirming that independence would have erased a material leverage channel.

\section{VaR, Expected Shortfall, and validation theory}

\subsection{Risk measures}

For scenario P\&L $\Pi$ and confidence $c$, with $\alpha=1-c$,
\begin{align}
\mathrm{VaR}_{c}&=-Q_{\alpha}(\Pi),\\
\mathrm{ES}_{c}&=-\mathbb{E}\!\left[\Pi\mid \Pi\le Q_{\alpha}(\Pi)\right].
\end{align}
A violation indicator is $I_t=\mathbf{1}\{\Pi_t<-\mathrm{VaR}_{c,t}\}$. A calibrated model should
have $\mathbb{E}[I_t]=\alpha$ and no serial dependence in $I_t$.

\subsection{Kupiec unconditional coverage}

For $x=\sum_t I_t$ exceptions in $n$ forecasts, $\widehat p=x/n$, the likelihood-ratio statistic is
\begin{equation}
LR_{UC}=-2\log\left[
\frac{(1-\alpha)^{n-x}\alpha^x}
{(1-\widehat p)^{n-x}\widehat p^x}
\right]\ \overset{H_0}{\sim}\ \chi^2_1.
\end{equation}
It checks the average exception frequency but cannot detect clustering.

\subsection{Christoffersen independence and conditional coverage}

Let $n_{ij}$ count transitions from $I_{t-1}=i$ to $I_t=j$, with
$\pi_0=n_{01}/(n_{00}+n_{01})$, $\pi_1=n_{11}/(n_{10}+n_{11})$, and
$\pi=(n_{01}+n_{11})/\sum n_{ij}$. Comparing the independent Bernoulli likelihood with the
first-order Markov likelihood gives $LR_{IND}\sim\chi^2_1$. The joint statistic is
\begin{equation}
LR_{CC}=LR_{UC}+LR_{IND}\ \overset{H_0}{\sim}\ \chi^2_2.
\end{equation}
This directly reproduces the logic of the legacy UC, IND, and CC MATLAB tests in maintained Python.

\subsection{Loss comparison and Murphy diagrams}

Exception counts ignore the economic size and placement of forecasts. The consistent lower-tail
quantile loss for forecast $q_t=-\mathrm{VaR}_{c,t}$ is
\begin{equation}
L_{\alpha}(q_t,y_t)=\left[\alpha-\mathbf{1}\{y_t<q_t\}\right](y_t-q_t).
\end{equation}
Pairwise Diebold-Mariano tests use the loss differential and a Newey-West long-run variance with 15
lags, chosen as $\lfloor n^{1/3}\rfloor$. Murphy diagrams integrate the elementary quantile score
\begin{equation}
S_{\alpha,\theta}(q,y)=
\left[\mathbf{1}\{y<q\}-\alpha\right]
\left[\mathbf{1}\{\theta<q\}-\mathbf{1}\{\theta<y\}\right]
\end{equation}
over thresholds $\theta$. A curve lying below another indicates better performance for that class of
consistent quantile scores; crossing curves reveal criterion-dependent rankings.

\section{Empirical VaR results at 90\%, 95\%, and 99\%}

\subsection{Coverage and formal tests}

\begin{table}[H]
\centering
\caption{Multi-level VaR performance and coverage tests. Pass/reject refers to the 5\% joint
conditional-coverage test.}
\label{tab:var-diagnostics}
\resizebox{\textwidth}{!}{%
\begin{tabular}{llrrrrrrc}
\toprule
Model & VaR & Breaches / expected & Rate & UC $p$ & IND $p$ & CC $p$ & Mean quantile loss & Decision \\
\midrule
PSP & 90\% & 386 / 402.7 & 9.59\% & 0.377 & 0.713 & 0.633 & 897.6 & Pass \\
PSP & 95\% & 186 / 201.4 & 4.62\% & 0.261 & 0.408 & 0.377 & 597.5 & Pass \\
PSP & 99\% & 41 / 40.3 & 1.02\% & 0.908 & 0.358 & 0.651 & 223.5 & Pass \\
PCA & 90\% & 402 / 402.7 & 9.98\% & 0.971 & 0.746 & 0.948 & 916.4 & Pass \\
PCA & 95\% & 213 / 201.4 & 5.29\% & 0.404 & 0.406 & 0.500 & 624.3 & Pass \\
PCA & 99\% & 48 / 40.3 & 1.19\% & 0.235 & 0.135 & 0.161 & 233.0 & Pass \\
PWG & 90\% & 255 / 402.7 & 6.33\% & $<0.001$ & 0.322 & $<0.001$ & 954.7 & Reject \\
PWG & 95\% & 155 / 201.4 & 3.85\% & $<0.001$ & 0.115 & $<0.001$ & 624.6 & Reject \\
PWG & 99\% & 62 / 40.3 & 1.54\% & 0.001 & 0.088 & 0.001 & 230.4 & Reject \\
\bottomrule

\end{tabular}}
\end{table}

PSP produces 386, 186, and 41 violations at 90\%, 95\%, and 99\%, respectively, against expectations
of 402.7, 201.4, and 40.3. PCA is exceptionally close at 90\% and remains statistically calibrated at
all levels. PWG produces too few exceptions at 90\% and 95\%, but too many at 99\%, rejecting correct
unconditional and conditional coverage throughout.

\begin{figure}[H]
\centering
\includegraphics[width=\textwidth]{15_var_coverage_summary.png}
\caption{Observed exception rates and 95\% Wilson confidence intervals against each nominal tail
probability.}
\label{fig:var-coverage}
\end{figure}

\begin{figure}[H]
\centering
\includegraphics[width=\textwidth]{16_var_test_heatmap.png}
\caption{Coverage-test $p$-values. The PWG unconditional and joint results are rejected at all three
levels, while PSP and PCA remain above 5\%.}
\label{fig:var-tests}
\end{figure}

\subsection{Loss, model ranking, and Diebold-Mariano tests}

\begin{table}[H]
\centering
\caption{Transparent ranking by nominal coverage accuracy and consistent quantile loss}
\label{tab:ranking}
\begin{tabular}{llrrrrr}
\toprule
VaR & Model & Absolute coverage error & Mean loss & Coverage rank & Loss rank & Overall rank \\
\midrule
90\% & PCA & 0.017\% & 916.4 & 1 & 2 & 1 \\
90\% & PSP & 0.415\% & 897.6 & 2 & 1 & 1 \\
90\% & PWG & 3.668\% & 954.7 & 3 & 3 & 3 \\
95\% & PCA & 0.289\% & 624.3 & 1 & 2 & 1 \\
95\% & PSP & 0.381\% & 597.5 & 2 & 1 & 1 \\
95\% & PWG & 1.151\% & 624.6 & 3 & 3 & 3 \\
99\% & PSP & 0.018\% & 223.5 & 1 & 1 & 1 \\
99\% & PCA & 0.192\% & 233.0 & 2 & 3 & 2 \\
99\% & PWG & 0.540\% & 230.4 & 3 & 2 & 2 \\
\bottomrule

\end{tabular}
\end{table}

The combined rank deliberately reports ties. At 90\% and 95\%, PCA ranks first on coverage accuracy
while PSP ranks first on quantile loss. PSP is the clear 99\% leader. This is more informative than a
single opaque composite score.

\begin{table}[H]
\centering
\caption{Pairwise Diebold-Mariano tests using HAC quantile-loss differentials. A negative mean
difference favors the first model.}
\label{tab:dm}
\begin{tabular}{llrrrl}
\toprule
VaR & Pair & DM statistic & $p$-value & Mean loss difference & 5\% interpretation \\
\midrule
90\% & PSP--PCA & -3.585 & $<0.001$ & -18.8 & PSP \\
90\% & PSP--PWG & -3.579 & $<0.001$ & -57.1 & PSP \\
90\% & PCA--PWG & -2.223 & 0.026 & -38.3 & PCA \\
95\% & PSP--PCA & -2.628 & 0.009 & -26.8 & PSP \\
95\% & PSP--PWG & -2.489 & 0.013 & -27.2 & PSP \\
95\% & PCA--PWG & -0.023 & 0.982 & -0.3 & No difference \\
99\% & PSP--PCA & -1.400 & 0.162 & -9.5 & No difference \\
99\% & PSP--PWG & -0.828 & 0.408 & -6.9 & No difference \\
99\% & PCA--PWG & 0.425 & 0.671 & 2.6 & No difference \\
\bottomrule

\end{tabular}
\end{table}

At 90\%, PSP significantly improves on PCA and PWG, and PCA significantly improves on PWG. At 95\%,
PSP is significantly better than both alternatives, whereas PCA and PWG have indistinguishable mean
loss. At 99\%, the sample supplies no statistically significant pairwise distinction despite PSP's
better coverage and average score.

\begin{figure}[H]
\centering
\includegraphics[width=\textwidth]{17_var_murphy_diagrams.png}
\caption{Murphy diagrams across all three VaR levels. Crossings make clear why coverage, average loss,
and pairwise significance should be reported together.}
\label{fig:murphy}
\end{figure}

\subsection{Correlation-PCA and Heston extensions}

\begin{table}[H]
\centering
\caption{Additional out-of-sample models under the same portfolio and validation tests}
\label{tab:extended-diagnostics}
\resizebox{\textwidth}{!}{%
\begin{tabular}{llrrrrrrl}
\toprule
Model & VaR & Breaches / expected & Rate & UC $p$ & IND $p$ & CC $p$ & Mean quantile loss & Decision \\
\midrule
PCA-C & 90\% & 410 / 402.7 & 10.18\% & 0.702 & 0.701 & 0.863 & 917.2 & Pass \\
PCA-C & 95\% & 205 / 201.4 & 5.09\% & 0.792 & 0.618 & 0.853 & 624.3 & Pass \\
PCA-C & 99\% & 46 / 40.3 & 1.14\% & 0.375 & 0.555 & 0.567 & 234.1 & Pass \\
Heston-HS & 90\% & 371 / 402.7 & 9.21\% & 0.092 & 0.735 & 0.228 & 896.6 & Pass \\
Heston-HS & 95\% & 186 / 201.4 & 4.62\% & 0.261 & 0.248 & 0.273 & 595.6 & Pass \\
Heston-HS & 99\% & 36 / 40.3 & 0.89\% & 0.491 & 0.420 & 0.570 & 218.8 & Pass \\
\bottomrule

\end{tabular}}
\end{table}

PCA-C preserves formal calibration at every level and puts the canonical-factor interpretation on a
transparent statistical footing. It has 410, 205, and 46 breaches. Daily calibrated Heston-HS has
\HestonNinetyBreaches, \HestonNinetyFiveBreaches, and \HestonNinetyNineBreaches\ breaches and also
passes UC, IND, and CC at all levels. Its average quantile loss is the lowest of the five methods at
90\%, 95\%, and 99\%.

\begin{figure}[H]
\centering
\includegraphics[width=\textwidth]{22_extended_var_coverage.png}
\caption{Observed coverage for the three original methods, correlation-PCA, and Heston-HS. Labels
show the 5\% conditional-coverage decision.}
\label{fig:extended-coverage}
\end{figure}

\begin{table}[H]
\centering
\caption{PSP minus Heston-HS HAC quantile-loss comparisons. Positive mean differences favor Heston.}
\label{tab:heston-dm}
\begin{tabular}{rrrr}
\toprule
VaR & DM statistic & $p$-value & Mean loss difference \\
\midrule
90\% & 0.629 & 0.529 & 1.0 \\
95\% & 1.187 & 0.235 & 1.8 \\
99\% & 1.436 & 0.151 & 4.7 \\
\bottomrule

\end{tabular}
\end{table}

Heston's loss improvement over PSP is small and not statistically significant at 5\%. The result is
therefore evidence that a daily stochastic-volatility projection is competitive, not evidence that
Heston dominates the nonparametric history.

\subsection{Heston surface fit and parameter diagnostics}

\begin{figure}[H]
\centering
\includegraphics[width=\textwidth]{20_heston_surface_3d.png}
\caption{Latest observed surface, daily calibrated Heston-implied surface, and error. Heston prices
are inverted to Black--Scholes implied volatility solely to retain the common grid and exposure map.}
\label{fig:heston-surface}
\end{figure}

Across all 4,280 dates, optimization reports success on every date; median daily IV RMSE is
\HestonMedianRMSE\ and the 90th percentile is \HestonPninetyRMSE. The following parameter ranges
show that a one-factor Heston surface is useful but strained by the data.

\begin{table}[H]
\centering
\caption{Daily Heston parameter distribution}
\label{tab:heston-parameters}
\begin{tabular}{lrrr}
\toprule
Parameter & 10th percentile & Median & 90th percentile \\
\midrule
$\sqrt{v_0}$ & 9.5\% & 15.1\% & 28.2\% \\
$\kappa$ & 1.294 & 3.591 & 7.892 \\
$\sqrt{\theta}$ & 16.1\% & 23.1\% & 35.9\% \\
$\xi$ & 0.560 & 1.151 & 1.500 \\
$\rho$ & -0.842 & -0.771 & -0.690 \\
\bottomrule

\end{tabular}
\end{table}

Only \HestonFellerShare\ of dates satisfy $2\kappa\theta\ge\xi^2$, and $\xi$ reaches its disclosed
upper bound on \HestonXiBoundShare\ of dates. These are material warnings about parameter
identification and single-factor Heston rigidity. They do not invalidate the common-grid risk
exercise, but they rule out treating every daily parameter estimate as a stable structural truth.

\begin{figure}[H]
\centering
\includegraphics[width=\textwidth]{21_heston_parameters.png}
\caption{Daily Heston parameters and surface-fit RMSE. Crisis spikes in $v_0$ are economically
interpretable; frequent $\xi$ boundary contacts document model tension rather than hiding it.}
\label{fig:heston-parameters}
\end{figure}

\subsection{Joint Heston simulation and full-revaluation evidence}

This subsection changes both the scenario law and the realized P\&L target. It therefore does not
pool ranks or dollar losses with the earlier vega-only experiment. All three methods below forecast
the same 4,027 realizations of the same fixed-strike, delta-hedged call strip.

\begin{table}[H]
\centering
\caption{Trailing physical Heston proxy; ``bound share'' records numerical stabilization rather
than deleting difficult dates}
\label{tab:physical-heston}
\begin{tabular}{lrrrr}
\toprule
Parameter & 10th percentile & Median & 90th percentile & Bound share \\
\midrule
$\mu_P$ & -0.056 & 0.112 & 0.213 & -- \\
$\kappa_P$ & 15.482 & 35.055 & 81.847 & 7.7\% \\
$\theta_P$ & 0.015 & 0.024 & 0.104 & 0.0\% \\
$\xi_P$ & 0.609 & 0.852 & 1.327 & 0.0\% \\
$\rho_P$ & -0.749 & -0.699 & -0.436 & 0.0\% \\
\bottomrule

\end{tabular}
\end{table}

The median $\rho_P=-0.699$ displays the expected leverage sign. Mean reversion is much less stable:
$\kappa_P$ has median 35.1 and contacts its cap on \PhysicalKappaBoundShare\ of dates. This is a
warning that daily option-implied $v_0$ is a noisy physical-state proxy, not evidence of literal
35-times annual structural mean reversion.

\begin{figure}[H]
\centering
\includegraphics[width=\textwidth]{25_heston_mc_diagnostics.png}
\caption{Joint Heston simulation diagnostics. The top panels show trailing physical parameters; the
bottom panels compare realized and simulated return tails and the 4/8/16-step discretization check.}
\label{fig:heston-mc-diagnostics}
\end{figure}

\begin{table}[H]
\centering
\caption{Full-revaluation VaR and ES for one common fixed-strike, delta-hedged portfolio. Dollar
columns are daily averages; pass/reject uses the 5\% joint conditional-coverage test.}
\label{tab:full-revaluation}
\resizebox{\textwidth}{!}{%
\begin{tabular}{llrrrrrrrrl}
\toprule
Model & VaR & Breaches / expected & Rate & Mean VaR & Mean ES & UC $p$ & IND $p$ & CC $p$ & Q-loss & Decision \\
\midrule
GBM+PSP-FHS & 90\% & 382 / 402.7 & 9.49\% & 2,798 & 5,124 & 0.273 & 0.038 & 0.064 & 455.0 & Pass \\
GBM+PSP-FHS & 95\% & 187 / 201.4 & 4.64\% & 4,075 & 6,884 & 0.294 & $<0.001$ & $<0.001$ & 307.2 & Reject \\
GBM+PSP-FHS & 99\% & 37 / 40.3 & 0.92\% & 8,470 & 12,455 & 0.600 & 0.004 & 0.015 & 125.7 & Reject \\
Heston-MC-P & 90\% & 333 / 402.7 & 8.27\% & 4,482 & 5,855 & $<0.001$ & 0.038 & $<0.001$ & 534.1 & Reject \\
Heston-MC-P & 95\% & 200 / 201.4 & 4.97\% & 5,567 & 6,722 & 0.922 & 0.751 & 0.947 & 331.7 & Pass \\
Heston-MC-P & 99\% & 72 / 40.3 & 1.79\% & 7,439 & 8,292 & $<0.001$ & 0.788 & $<0.001$ & 100.2 & Reject \\
Heston-MC-Q & 90\% & 110 / 402.7 & 2.73\% & 4,342 & 5,553 & $<0.001$ & $<0.001$ & $<0.001$ & 515.9 & Reject \\
Heston-MC-Q & 95\% & 73 / 201.4 & 1.81\% & 5,307 & 6,312 & $<0.001$ & 0.002 & $<0.001$ & 330.9 & Reject \\
Heston-MC-Q & 99\% & 45 / 40.3 & 1.12\% & 6,932 & 7,669 & 0.462 & $<0.001$ & $<0.001$ & 118.7 & Reject \\
\bottomrule

\end{tabular}}
\end{table}

The paired GBM+PSP-FHS benchmark records \FullGBMNinetyBreaches,
\FullGBMNinetyFiveBreaches, and \FullGBMNinetyNineBreaches\ exceptions. Its unconditional rates are
close to nominal at every level. Joint coverage passes at 90\% ($p=0.064$), although the separate
independence test is marginally rejected, and fails at 95\% and 99\% because exceptions cluster.

Heston-MC-P records \FullHestonPNinetyBreaches, \FullHestonPNinetyFiveBreaches, and
\FullHestonPNinetyNineBreaches\ exceptions. It is conservative at 90\%, essentially exact at 95\%
(200 versus 201.4 expected, CC $p=\FullHestonPNinetyFiveCCP$), and anti-conservative at 99\% (72
versus 40.3). That pattern is economically important: the diffusion captures an intermediate tail
but not the most extreme discontinuous or regime-changing losses. Heston-MC-Q fails joint coverage
at all levels. Its near-nominal 99\% unconditional count is not a rescue because its exceptions are
strongly dependent; more fundamentally, $Q$ is not the physical forecasting law.

\begin{figure}[H]
\centering
\includegraphics[width=\textwidth]{26_full_revaluation_coverage.png}
\caption{Observed full-revaluation breach rates. Labels report the joint conditional-coverage
decision, not only proximity to the dashed nominal rate.}
\label{fig:full-revaluation-coverage}
\end{figure}

\begin{table}[H]
\centering
\caption{HAC Diebold--Mariano comparisons for full-revaluation quantile loss. The loss difference is
method A minus method B.}
\label{tab:full-revaluation-dm}
\resizebox{\textwidth}{!}{%
\begin{tabular}{llrrrl}
\toprule
VaR & Comparison & DM statistic & $p$-value & Mean loss difference & Preferred at 5\% \\
\midrule
90\% & GBM+PSP-FHS--Heston-MC-P & -4.812 & $<0.001$ & -79.2 & GBM+PSP-FHS \\
90\% & GBM+PSP-FHS--Heston-MC-Q & -6.015 & $<0.001$ & -60.9 & GBM+PSP-FHS \\
90\% & Heston-MC-P--Heston-MC-Q & 1.269 & 0.204 & 18.3 & No difference \\
95\% & GBM+PSP-FHS--Heston-MC-P & -1.790 & 0.073 & -24.6 & No difference \\
95\% & GBM+PSP-FHS--Heston-MC-Q & -3.195 & 0.001 & -23.8 & GBM+PSP-FHS \\
95\% & Heston-MC-P--Heston-MC-Q & 0.062 & 0.950 & 0.8 & No difference \\
99\% & GBM+PSP-FHS--Heston-MC-P & 3.606 & $<0.001$ & 25.5 & Heston-MC-P \\
99\% & GBM+PSP-FHS--Heston-MC-Q & 0.734 & 0.463 & 7.0 & No difference \\
99\% & Heston-MC-P--Heston-MC-Q & -1.440 & 0.150 & -18.5 & No difference \\
\bottomrule

\end{tabular}}
\end{table}

Loss rankings reinforce that no engine dominates. GBM+PSP-FHS has significantly lower 90\% loss
than Heston-MC-P; the 95\% difference is not significant; and Heston-MC-P has significantly lower
99\% loss despite its excessive exception count. Quantile loss rewards threshold placement and
exception severity, so it need not rank models in the same order as coverage.

\begin{figure}[p]
\centering
\includegraphics[width=\textwidth]{27_full_revaluation_backtest_95.png}
\caption{Realized nonlinear P\&L and 95\% thresholds. Large positive P\&L spikes are retained so the
scale of the complete revaluation target is visible; VaR violations occur below the colored lines.}
\label{fig:full-revaluation-paths}
\end{figure}

\begin{figure}[p]
\centering
\includegraphics[width=\textwidth]{28_full_revaluation_regimes.png}
\caption{Observed-to-expected breach ratios by regime. Heston-MC-P is conservative in the GFC at
90\%--95\%, close at GFC 99\%, anti-conservative in the expansion 99\% tail, and strongly conservative
during COVID.}
\label{fig:full-revaluation-regimes}
\end{figure}
\FloatBarrier

\begin{table}[H]
\centering
\caption{Projected-Euler sensitivity at twelve representative dates}
\label{tab:heston-numerical}
\begin{tabular}{lrrrr}
\toprule
Measure & Substeps & Median simulated 1\% return & Max mean error & Max zero-variance share \\
\midrule
Heston-P & 4 & -2.368\% & 0.033\% & 0.100\% \\
Heston-P & 8 & -2.374\% & 0.033\% & 0.040\% \\
Heston-P & 16 & -2.388\% & 0.034\% & 0.020\% \\
Heston-Q & 4 & -2.545\% & 0.043\% & 2.880\% \\
Heston-Q & 8 & -2.608\% & 0.043\% & 1.740\% \\
Heston-Q & 16 & -2.632\% & 0.044\% & 1.000\% \\
\bottomrule

\end{tabular}
\end{table}

The common-random-number 4/8/16-step check shows small changes relative to daily tail scale; across
the sampled dates the largest absolute P-measure mean-return error is \HestonMCMeanError. The
bilinear conditional-price grid differs from direct Heston inversion by at most
\HestonPriceGridError\ basis points of normalized spot. The production eight-step scheme can place a
small fraction of paths exactly at zero variance, disclosed in \cref{tab:heston-numerical}; exact
affine simulation would be a higher-cost robustness extension \cite{broadiekaya2006}.

At realized repricing, \FullExtrapolatedShare\ of node evaluations lie just outside the original
fixed grid, primarily because the 30-day bucket rolls to 29 days and boundary strikes move in
moneyness. These values use local edge extrapolation and are explicitly counted. Only
\CurveFallbackCount\ forecast origin lacks a same-day zero curve; it uses the most recent prior WRDS
curve (three calendar days old) without looking ahead.

\subsection{Specification, crisis, and stability robustness}

Five deterministic PSP variants combine 250-, 500-, and 750-day windows with decay 0.985, decay
0.995, or equal weighting. None rejects joint conditional coverage at any confidence level; observed
to expected breach ratios range from 0.87 to 1.17. Subperiod tests cover the GFC, 2010--2019
expansion, and COVID period. PCA, PCA-C, and Heston-HS pass CC in all twelve full/subperiod-level
tests; PSP fails only COVID 95\% independence, whereas PWG fails nine of twelve.

\begin{figure}[H]
\centering
\includegraphics[width=\textwidth]{23_robustness_dashboard.png}
\caption{Left: PSP sensitivity to rolling window and decay. Right: 99\% full-sample and regime
coverage. A ratio of one is nominal.}
\label{fig:robustness}
\end{figure}

\begin{figure}[H]
\centering
\includegraphics[width=\textwidth]{24_extended_var_backtest_95.png}
\caption{Observed daily vega P\&L and 95\% VaR thresholds for all five methods.}
\label{fig:extended-paths}
\end{figure}
\FloatBarrier

\subsection{Forecast paths and violations}

The panels below expose the timing and clustering of exceptions, complementing the aggregate UC,
IND, and CC statistics. The same realized P\&L is used at every level so the comparison isolates the
forecast threshold rather than a change in portfolio or evaluation sample.

\begin{figure}[H]
\centering
\includegraphics[width=.97\textwidth]{13_var_backtest_90.png}
\caption{Out-of-sample 90\% VaR paths and exceptions for PSP, PCA, and PWG.}
\label{fig:var90}
\end{figure}

\begin{figure}[H]
\centering
\includegraphics[width=.97\textwidth]{06_var_backtest_95.png}
\caption{Out-of-sample 95\% VaR paths and exceptions for PSP, PCA, and PWG.}
\label{fig:var95}
\end{figure}

\begin{figure}[H]
\centering
\includegraphics[width=.97\textwidth]{14_var_backtest_99.png}
\caption{Out-of-sample 99\% VaR paths and exceptions for PSP, PCA, and PWG.}
\label{fig:var99}
\end{figure}
\FloatBarrier

\begin{figure}[H]
\centering
\includegraphics[width=.82\textwidth]{07_pnl_distribution.png}
\caption{Empirical distribution of realized daily vega P\&L. The visibly non-Gaussian tail motivates
historical and factor-based scenarios.}
\label{fig:pnl}
\end{figure}

\section{Computational and reproducibility design}

\subsection{Efficiency measures}

\begin{itemize}
  \item Annual WRDS tables are queried with server-side date, ticker, option-type, price, maturity,
  and moneyness restrictions; completed years are cached as Parquet and not downloaded again.
  \item Surface partitions are keyed by a configuration fingerprint. A changed filter or surface
  grid cannot silently reuse an incompatible cache.
  \item Only the compact 45-node daily surface is combined across 17 years, rather than retaining
  millions of option rows in memory.
  \item Randomized rolling PCA, vectorized vega P\&L, scalar Gaussian portfolio draws, and a fixed
  random seed make the risk stage efficient and reproducible.
  \item Heston uses vectorized strikes by maturity, damped Fourier inversion, cached
  Gauss--Laguerre nodes, warm-started daily optimization, and a cache signature tied to every Heston
  setting and grid coordinate.
  \item Full revaluation vectorizes 2,000 state paths, evaluates conditional Heston prices on compact
  variance--strike grids, interpolates by scenario, and fingerprints its Monte Carlo cache separately
  from the daily surface-calibration cache.
  \item The multi-level VaR report is refreshed from the saved backtest with
  \path{scripts/refresh_var_evaluation.py}; no WRDS re-download is needed.
\end{itemize}

\subsection{Research safeguards}

\begin{itemize}
  \item Forecast windows include data dated through origin $t$ and exclude the $t+1$ realization,
  eliminating look-ahead from scenario generation.
  \item Surface method comparison is held out by date, not evaluated on fitting residuals alone.
  \item Unsupported grid nodes are explicitly identified and stabilized with a nearest-fit hull
  guard; \ExtrapolatedShare\ of the long-run grid nodes carry that flag.
  \item The exact configuration fixes the time interval, surface grid, random seed, risk window,
  scenario count, and confidence levels.
  \item WRDS credentials stay in a mode-600, Git-ignored \code{.env}. No password, phone number,
  SMS code, or MFA artifact is written to source, output metadata, the LaTeX report, or Git.
\end{itemize}

\subsection{Reproduction commands}

\begin{verbatim}
make install
make test
make demo
make wrds

# Refresh 90%, 95%, 99% diagnostics and figures from saved forecasts
PYTHONPATH=src python scripts/refresh_var_evaluation.py \
  --config configs/long_horizon.yaml

# Reproduce PCA diagnostics, Heston calibration, joint MC, and robustness
PYTHONPATH=src python scripts/run_extended_analysis.py \
  --config configs/long_horizon.yaml

# Regenerate report tables, then compile the LaTeX source
python reports/latex/generate_tables.py
latexmk -pdf reports/latex/SPX_WRDS_Thesis_Report_2005_2021.tex
\end{verbatim}

\section{Limitations and extensions}

\begin{enumerate}
  \item \textbf{Two targets, not one horse race.} The vega-only experiment isolates surface risk; the
  full-revaluation experiment adds spot, convexity, carry, and model-pricing effects. Their dollar
  losses and rankings answer different questions and must not be pooled.
  \item \textbf{Synthetic exposure shape.} Both experiments use the disclosed 500,000-vega profile,
  not a live portfolio. Full revaluation uses fixed strikes and economically scaled contract
  quantities, but transaction costs, bid--ask execution, discrete hedge trading, and portfolio
  import remain outside scope.
  \item \textbf{Surface arbitrage.} Positivity and hull controls do not impose calendar and butterfly
  no-arbitrage constraints. An SVI/SSVI or constrained spline benchmark would strengthen the surface.
  \item \textbf{Expected Shortfall validation.} ES paths are produced, but the report emphasizes VaR
  because the legacy tests are VaR-oriented. Joint VaR-ES scoring and Acerbi-Szekely style tests are
  natural additions. At 99\%, 2,000 scenarios leave only about 20 draws in the ES tail, so ES is less
  statistically precise than VaR.
  \item \textbf{Regime dependence.} The 500-day rolling window adapts gradually. Explicit volatility
  regimes or filtered historical simulation may react faster during structural breaks.
  \item \textbf{Heston identification.} Only \HestonFellerShare\ of daily fits satisfy the Feller
  inequality and $\xi$ reaches its cap on \HestonXiBoundShare. Multi-start calibration,
  piecewise-constant parameters, rough Heston, or jumps could reduce one-factor misspecification,
  but they would add parameters and estimation risk.
  \item \textbf{Physical-state proxy and repricing freeze.} Heston-MC-P treats calibrated $v_0$ as an
  observable proxy for physical variance; \PhysicalKappaBoundShare\ of $\kappa_P$ estimates require
  stabilization. Horizon pricing freezes the current $Q$ structural parameters and the additive
  market-price residual for one day. Joint P/Q estimation or a latent-state filter would be stronger.
  \item \textbf{Diffusion and discretization.} A continuous Heston diffusion cannot generate price
  jumps and the production scheme is projected Euler. The 4/8/16 check is reassuring for one-day
  tails but does not replace exact simulation, quadratic-exponential discretization, or jump-Heston
  robustness.
\end{enumerate}

\section{Conclusion}

The reorganized project now exposes the full chain from WRDS data to surface geometry, joint state
simulation, nonlinear repricing, and formal tail-risk evaluation. In the controlled surface-shock
experiment, PSP, covariance-PCA, correlation-PCA, and Heston-HS are calibrated at all three levels;
Heston-HS has the lowest average quantile loss but does not significantly dominate PSP. The Gaussian
PWG benchmark demonstrates why shrinkage alone does not repair a misspecified tail distribution.

The new full-revaluation experiment validates the user's intuition: if Heston is proposed as a market
dynamics model, both future spot and variance should be simulated. Doing so produces a more demanding
answer than surface reconstruction alone. Heston-MC-P is exceptionally well calibrated at 95\%, but
fails at 90\% and underestimates the 99\% tail; Heston-MC-Q fails as a physical VaR law. The paired
GBM+PSP-FHS benchmark obtains the closest unconditional rates but still exhibits exception clustering.
The evidence supports Heston as a useful structured scenario model, not as a universally adequate
tail law.

The PCA correction is equally important. Raw covariance PCA is useful for physical-IV reconstruction,
but it does not justify automatically naming PC1--PC3 level, slope, and curvature. Correlation scaling
recovers level and moneyness slope, while the two-dimensional surface assigns PC3 to term slope and
PC5 to smile curvature. Surface reconstruction, factor interpretation, scenario generation, and VaR
validation must be assessed separately, then interpreted jointly.

\appendix
\section{Additional diagnostic figures}

\begin{figure}[H]
\centering
\includegraphics[width=.92\textwidth]{03_regression_quality.png}
\caption{Daily production-surface fit and cross-sectional log-IV residual distribution.}
\label{fig:regression-appendix}
\end{figure}

\begin{figure}[H]
\centering
\includegraphics[width=\textwidth]{05_pca_components.png}
\caption{Heat-map representation of the leading PCA loading surfaces, complementing
\cref{fig:pca-3d}.}
\label{fig:pca-heatmap-appendix}
\end{figure}

\section{Result and code artifact map}

\begin{longtable}{p{.38\textwidth}p{.54\textwidth}}
\caption{Reproducible artifact map}\label{tab:artifacts}\\
\toprule
Artifact & Purpose \\
\midrule
\endfirsthead
\toprule
Artifact & Purpose \\
\midrule
\endhead
\path{configs/long_horizon.yaml} & Exact 2005--2021 data, surface, PCA, and risk configuration. \\
\path{src/spx_risk/data/wrds.py} & WRDS discovery, annual server-side extraction, normalization, and caching. \\
\path{analysis/black_scholes.py} & European pricing, Greeks, and Brent implied-volatility inversion. \\
\path{analysis/surface.py} & Filters, forward moneyness, four surface estimators, holdout tests, and hull guard. \\
\path{analysis/pca.py} & Fixed-grid surface matrix, daily changes, PCA factors, loadings, and reconstruction. \\
\path{analysis/risk.py} & PSP/PCA/PWG scenarios, VaR/ES, UC/IND/CC tests, loss comparison, and rankings. \\
\path{analysis/heston.py} & Heston Fourier prices, IV inversion, daily calibration, diagnostics, and model surfaces. \\
\path{analysis/heston\_mc.py} & Physical-parameter proxy, paired FHS, joint Heston paths, fixed-strike full revaluation, and numerical checks. \\
\path{analysis/robustness.py} & Canonical PCA tests, subspace stability, PSP sensitivity, Heston-HS, and subperiod tests. \\
\path{visualization/plots.py} & Publication figures, 3D surfaces, multi-level backtests, coverage heat map, and Murphy diagrams. \\
\path{visualization/extended.py} & PCA comparison, Heston surfaces and parameters, extended coverage, and robustness figures. \\
\path{tables/var_backtest.csv} & 4,279 dated P\&L rows with VaR, ES, and breaches for nine model-level combinations. \\
\path{tables/backtest_diagnostics.csv} & Exception counts, severity, loss, UC, IND, and CC statistics. \\
\path{tables/diebold_mariano_comparison.csv} & Nine HAC pairwise quantile-loss comparisons. \\
\path{tables/var_model_ranking.csv} & Coverage and quantile-loss ranks at 90\%, 95\%, and 99\%. \\
\path{tables/extended_backtest_diagnostics.csv} & Five-model VaR diagnostics including PCA-C and Heston-HS. \\
\path{tables/pca_interpretation_diagnostics.csv} & Canonical shape cosines for covariance and correlation PCA. \\
\path{tables/heston_calibration_diagnostics.csv} & Daily optimizer, IV-error, Feller, and parameter audit trail. \\
\path{tables/full_revaluation_diagnostics.csv} & Three-engine nonlinear VaR/ES, exception, loss, and UC/IND/CC results. \\
\path{tables/heston_mc_numerical_robustness.csv} & P/Q step convergence, mean-return errors, boundary mass, and price-grid accuracy. \\
\path{reports/latex/} & LaTeX source, generated result fragments, and this compiled report. \\
\bottomrule
\end{longtable}

\section{References}

\begin{thebibliography}{99}

\bibitem{blackscholes1973}
F. Black and M. Scholes.
\newblock The pricing of options and corporate liabilities.
\newblock \emph{Journal of Political Economy}, 81(3):637--654, 1973.

\bibitem{merton1973}
R. C. Merton.
\newblock Theory of rational option pricing.
\newblock \emph{Bell Journal of Economics and Management Science}, 4(1):141--183, 1973.

\bibitem{heston1993}
S. L. Heston.
\newblock A closed-form solution for options with stochastic volatility with applications to bond
and currency options.
\newblock \emph{Review of Financial Studies}, 6(2):327--343, 1993.

\bibitem{carrmadan1999}
P. Carr and D. Madan.
\newblock Option valuation using the fast Fourier transform.
\newblock \emph{Journal of Computational Finance}, 2(4):61--73, 1999.

\bibitem{lord2010}
R. Lord, R. Koekkoek, and D. van Dijk.
\newblock A comparison of biased simulation schemes for stochastic volatility models.
\newblock \emph{Quantitative Finance}, 10(2):177--194, 2010.

\bibitem{broadiekaya2006}
M. Broadie and O. Kaya.
\newblock Exact simulation of stochastic volatility and other affine jump diffusion processes.
\newblock \emph{Operations Research}, 54(2):217--231, 2006.

\bibitem{dumas1998}
B. Dumas, J. Fleming, and R. E. Whaley.
\newblock Implied volatility functions: empirical tests.
\newblock \emph{Journal of Finance}, 53(6):2059--2106, 1998.

\bibitem{kupiec1995}
P. H. Kupiec.
\newblock Techniques for verifying the accuracy of risk measurement models.
\newblock \emph{Journal of Derivatives}, 3(2):73--84, 1995.

\bibitem{christoffersen1998}
P. F. Christoffersen.
\newblock Evaluating interval forecasts.
\newblock \emph{International Economic Review}, 39(4):841--862, 1998.

\bibitem{dm1995}
F. X. Diebold and R. S. Mariano.
\newblock Comparing predictive accuracy.
\newblock \emph{Journal of Business and Economic Statistics}, 13(3):253--263, 1995.

\bibitem{ledoitwolf2004}
O. Ledoit and M. Wolf.
\newblock A well-conditioned estimator for large-dimensional covariance matrices.
\newblock \emph{Journal of Multivariate Analysis}, 88(2):365--411, 2004.

\bibitem{ehm2016}
W. Ehm, T. Gneiting, A. Jordan, and F. Krueger.
\newblock Of quantiles and expectiles: consistent scoring functions, Choquet representations and
forecast rankings.
\newblock \emph{Journal of the Royal Statistical Society: Series B}, 78(3):505--562, 2016.

\bibitem{zamani2023}
S. Zamani et al.
\newblock Temporal volatility surface projection: parametric surface projection method for
derivatives portfolio risk management.
\newblock \emph{arXiv:2311.14985}, 2023.

\end{thebibliography}

\end{document}
